% ===================================================================
%  14. TAULA — Kalea. — Merkatariak.
%  Traduction 2026 d'apres texto/io, controlee sur texto/fr, texto/en
%  et texto/es. Consigne : voir texto/eu/10-taula-01.tex.
%
%  LE TABLEAU QUI AVAIT COUTE DEUX INVERSIONS AU SUEDOIS — le jet (49)
%  de la fontaine (48) et l'imperiale (59) de l'omnibus (58) — en coute
%  deux au basque, et pour la meme raison de fond, le genitif prepose.
%  Ecrites d'avance en apposition postposee, elles ne coutent rien.
% ===================================================================

%%K t14-apar-1 apar
\VUsaut{4.0mm}
\VUtitre{15.0pt}{60}{14. TAULA}
\VUsaut{3.0mm}

%%K t14-tit-1 sub
\VUpk{11.4pt}{40}{Kalea. --- Merkatariak.}
\VUsaut{3.0mm}

%%K t14-01-1 p
1. --- Nire adiskide Joanes, herrian bizi dena, hiru egun igarotzera etorri da
gureenera. \VUgras{Orain} arte ez zaio behin ere gertatu hiri handi bat ikustea,
eta hala egun osoak ematen ditugu paseatzen, eta nik ez dakiena azaltzen diot.

%%K t14-02-1 p
2. --- Lehenik \VUgras{behatokia} \textsuperscript{(1)} ikustera joan ginen, eta
asko interesatu zitzaion. \VUgras{kaletik} \textsuperscript{(3)} itzultzean --- non
\VUgras{turkiar bainuak} \textsuperscript{(2)} baitaude ---, \VUgras{ehorzketa} bat
\textsuperscript{(4)} topatu genuen. \VUgras{Hileta-segizioaren} buruan
\VUgras{apaizeria} zihoan, \VUgras{hilkutxaren} aurretik, zeinaren gainean
\VUgras{zerraldoa} jarri baitzuten. Lagun asko zihoan familiaren atzetik,
\VUgras{doluz}.

%%K t14-02-2 p
Segizioa igaro zenean, kalea zeharkatu genuen \VUgras{garagardotegi} handi baten
aurrean \textsuperscript{(5)}, non \VUgras{bezero} askok \VUgras{garagardoa}
edaten baitzuen terrazan, \VUgras{beirazko estalkiaren} \textsuperscript{(6)}
azpian.

%%K t14-03-1 p
3. --- \VUgras{Armarriak}, \VUgras{ardo eta pattar saltzailearen}
\textsuperscript{(7)} dendaren eta \VUgras{partitura-saltzailearen}
\textsuperscript{(10)} dendaren artean jarria, harritu zuen Joanes. Zer esan nahi
zuen galdetu zidan; esan nion Ingalaterrako \VUgras{kontsulatua}
\textsuperscript{(8)} zela, eta \VUgras{kontsula} \textsuperscript{(9)} erakutsi
nion, balkoian zegoena.

%%K t14-tit-2 sub
\VUpk{10.2pt}{40}{Erakusleihoei begira.}
\VUsaut{3.0mm}

%%K t14-04-1 p
4. --- Jakina, \VUgras{musika-tresnen}, \VUgras{harmoniumen},
\VUgras{organoen} eta \VUgras{pianoen} erakusketaren aurrean gelditu ginen;
\VUgras{partituren} eta piano-\VUgras{piezen} azal ilustratuak begiratu genituen
eta iragarki gisa zerbitzatzen zuen zurezko \VUgras{lira} urreztatua miretsi.

%%K t14-05-1 p
5. --- \VUgras{Kale-garbitzaileek} \textsuperscript{(11)} eta \VUgras{eskuila-makinak}
\textsuperscript{(12)} altxatutako hauts-lainoek azkarrago igarotzera behartu
gintuzten \VUgras{altzarigilearen} \textsuperscript{(13)} \VUgras{altzari-multzo}
ederren ondotik eta \VUgras{burdindegiko} \textsuperscript{(14)}
\VUgras{berogailu} eta \VUgras{lanparen} erakusketaren ondotik.

%%K t14-06-1 p
6. --- Leku horretan bertan espaloitik jaitsi behar izan genuen benetan, fardelez
beteta baitzegoen, \VUgras{altzari-furgoitik} \textsuperscript{(16)} jaisten
zituztenak, alokatu berri zuten \VUgras{dendara} \textsuperscript{(15)} sartzeko.

%%K t14-06-2 p
Horrek eragotzi zigun \VUgras{kotxegilearenean} \textsuperscript{(17)} zeuden
kotxe ederrak ikustea, baina ez zigun eragotzi gelditzea eta begiratzea nola
\VUgras{paketatzaileak} \textsuperscript{(19)} kaxa bat josten zuen bere
auzokoarentzat, \VUgras{bizikleta eta automobil saltzailearentzat}
\textsuperscript{(18)}. Gero, nahiko berandu zenez, etxera itzuli ginen.

%%K t14-tit-3 sub
\VUpk{10.2pt}{40}{Hirian barrena paseatzea.}
\VUsaut{3.0mm}

%%K t14-07-1 p
7. --- Hurrengo egunean amak proposatu zuen Joanesi argazkia ateratzea, eta
\VUgras{argazkilariarengana} \textsuperscript{(20)} eramateko eskatu zidan.
Bazkalondoren artistaren \VUgras{estudiora} joan ginen, non nahiko itxaron behar
izan baikenuen. Denbora igarotzeko, horman zintzilik zeuden \VUgras{erretratuak}
\textsuperscript{(22)} begiratzen genituen.

%%K t14-07-2 p
Gure txanda iritsi zenean, kontua laburra izan zen, eta nire adiskideak
\VUgras{makinaren} \textsuperscript{(21)} aurrean posatzen amaitu orduko, jaitsi
ginen.

%%K t14-08-1 p
8. --- Lehen solairuan \VUgras{ile-apaindegiaren} \textsuperscript{(23)} ate
irekitik ikusi genuen nola morroiak bezero bati ilea mozten zion.

%%K t14-08-2 p
Berriro kalera irtenda, \VUgras{tabako-dendaren} \textsuperscript{(24)} ondotik
igaro ginen. \VUgras{Farola} gorriak \textsuperscript{(25)} eta \VUgras{iragarki}
\textsuperscript{(26)} gisa zerbitzatzen zuen \VUgras{pipa handiak} hain
dibertitu zuten Joanes, non \VUgras{tabako-hautsezko}
\textsuperscript{(27)} usain-aldi baten inguruan solasean ari ziren bi agureren
kontra jo baitzuen.

%%K t14-tit-4 sub
\VUpk{10.2pt}{40}{Gozotegia.}
\VUsaut{3.0mm}

%%K t14-09-1 p
9. --- \VUgras{Diru-trukatzailearen} \textsuperscript{(29)}
\VUgras{erakusleihoan} ikusgai zeuden herrialde guztietako \VUgras{billeteek} eta
\VUgras{txanponek} une batez geldiarazi gintuzten, baina askariaren ordua zenez,
gehiago berandutu gabe \VUgras{gozotegian} \textsuperscript{(31)} sartu ginen.

%%K t14-10-1 p
10. --- Dendan zeuden \VUgras{gozoki} guztiak ikusita --- \VUgras{pastelak,
opilak, gailetak} eta \VUgras{goxoki} mota guztiak --- nekez aukeratu ahal izan
genuen. Azkenean Joanes \VUgras{krema-pastelen} alde geratu zen, eta nik
\VUgras{gerezi-tarta} txiki bat baino ez nuen hartu, gozoak \VUgras{hortzetako
mina} ematen baitit eta ez baitut sarriegi \VUgras{hortz-medikuarengana}
\textsuperscript{(30)} joan nahi.

%%K t14-tit-5 sub
\VUpk{10.2pt}{40}{Makinaz idaztea.}
\VUsaut{3.0mm}

%%K t14-11-1 p
11. --- Adiskideari \VUgras{idazmakina} erakusteko --- zeinaz entzun baitzuen
baina ikusi ez ---, \VUgras{bulego} batera \textsuperscript{(32)} sartu ginen, nire
\VUgras{lehengusuarenera} \textsuperscript{(34)}. Bere \VUgras{idazkariari}
\textsuperscript{(35)} gutunak diktatzen aurkitu genuen, zeina
\VUgras{takigrafoa} baita. Gutuna amaitu orduko, \VUgras{kopiatzaileak}
\textsuperscript{(33)} bere makinan kopiatzen zuen. Senidearen lana ez oztopatzeko,
minutu batzuk baino ez ginen egon haren etxean.

%%K t14-12-1 p
12. --- Lehengusuaren bulegoaren azpian \VUgras{erlojugile eta bitxigile} handi
bat \textsuperscript{(37)} bizi da, zeina gainera zilargilea baita. Bere denda
bikainean \VUgras{horma-erloju, tximiniako erloju, sakelako erloju} asko dago,
\VUgras{kate} eta \VUgras{bitxi} garestiak: \VUgras{perla-lepokoak}, urrezko
\VUgras{eskumuturrekoak}, \VUgras{belarritakoak} eta \VUgras{eraztunak}
\VUgras{harribitxiekin} eta \VUgras{diamanteekin}. Erakusleiho bat oso-osorik
\VUgras{zilarrari} eskainita dago eta zilarrezko \VUgras{mahai-tresnen} bilduma
handia hartzen du.

%%K t14-13-1 p
13. --- Kalearen izkina inguratzean, ohartu nintzen etxeko \VUgras{zenbakiaren}
ondoko \VUgras{plaka} \textsuperscript{(36)} aldatu zutela. Ozen esatera nindoan,
banda militarraren soinu alaiak entzun zirenean. Erregimentura korrika abiatu
ginen, pentsatu gabe \VUgras{kapelagilearen} \textsuperscript{(39)} eta
\VUgras{eskularrugilearen} \textsuperscript{(40)} denden ondotik igaroko zela eta
hain ondo kokatuta egongo ginela haren igarotzea ikusteko \VUgras{optikariaren}
\textsuperscript{(38)} erakusleihoaren aurrean geratu bagina, \VUgras{pintza-betaurrekoz,
betaurrekoz} eta \VUgras{prismatikoz} betea.

%%K t14-tit-6 sub
\VUpk{10.2pt}{40}{Erregimentuaren igarotzea.}
\VUsaut{3.0mm}

%%K t14-14-1 p
14. --- \VUgras{Erregimentuaren} \textsuperscript{(41)} buruan
\VUgras{danbor-nagusia} \textsuperscript{(42)} zihoan, haren atzetik
\VUgras{danborrariak} \textsuperscript{(43)}, \VUgras{turutariak}
\textsuperscript{(44)} eta \VUgras{banda} osoa. \VUgras{Jenerala}
\textsuperscript{(45)}, zeina plazan baitzegoen bere laguntzailearekin,
\VUgras{ur-zorrotadaren} \textsuperscript{(49)} ondoan gelditu zen ---
\VUgras{iturrikoaren} \textsuperscript{(48)} ondoan --- \VUgras{desfilea}
ikusteko.

%%K t14-14-2 p
\VUgras{Estatu Nagusiko ofizialek} \textsuperscript{(46)},
\VUgras{koronelak} eta \VUgras{teniente-koronelak} ezpataz agurtu zuten.
\VUgras{Komandanteak} beren \VUgras{batailoien} \textsuperscript{(47)} buruan
zihoazen, eta \VUgras{kapitain} bakoitzak bere \VUgras{konpainia} agintzen zuen,
eta \VUgras{tenienteek}, \VUgras{alferezek} eta \VUgras{sarjentuek} beren ohiko
lekuak hartzen zituzten beren jendearen ondoan.

%%K t14-15-1 p
15. --- Bidaiari asko bildu zen \VUgras{bidegurutzean} \textsuperscript{(50)};
dendariak --- \VUgras{aterki-saltzailea} \textsuperscript{(51)},
\VUgras{tindatzailea} \textsuperscript{(53)}, \VUgras{arma-saltzailea}
\textsuperscript{(54)} --- \VUgras{soldaduak} ikustera atera ziren. Kotxe guztiak
--- \VUgras{alokairuko kotxeak} \textsuperscript{(52)}, \VUgras{kotxe
partikularrak} \textsuperscript{(57)}, eta baita \VUgras{ureztatzeko upela}
\textsuperscript{(56)} ere --- espaloira estutu ziren.

%%K t14-tit-7 sub
\VUpk{10.2pt}{40}{Kaleko eszenak.}
\VUsaut{3.0mm}

%%K t14-15-2 p
\VUgras{Bidaiari-agenteak} \textsuperscript{(55)}, bere \VUgras{lagin-maletak}
eskuan zeramatzanak, bizkor zeharkatu zuen kalea, eta \VUgras{inperialean}
\textsuperscript{(59)} eserita zeudenak --- \VUgras{omnibusekoan}
\textsuperscript{(58)} --- itzultzen ziren ikuskizuna miresteko. \VUgras{Kale-kantari}
gaixoak \textsuperscript{(60)} bere \VUgras{eskuorganoarekin}
\textsuperscript{(61)} bere \VUgras{kanta} eten behar izan zuen, danborren
zaratak haren ahotsa itotzen baitzuen.

%%K t14-16-1 p
16. --- Erregimentua \VUgras{oihal-dendaren} \textsuperscript{(64)} parera iritsi
zenean, \VUgras{jostunek} \textsuperscript{(63)} eta \VUgras{jantzigileek}
\textsuperscript{(62)} beren \VUgras{josteko makinak} eta lana utzi zituzten
leihotik begiratzeko. \VUgras{Dendariak} \textsuperscript{(65)}, zeinak
\VUgras{jantzi} berriak \textsuperscript{(66)} jarri berri baitzituen bere
\VUgras{erakusleihoan}, \VUgras{oihal} \textsuperscript{(67)} eta
\VUgras{mihise} pieza guztiak jaso behar izan zituen --- espaloian jarrita
zeuden \VUgras{ur-hobiaren} \textsuperscript{(68)} ondoan ---, jendetzak
irauliko zituen beldurrez; \VUgras{saltzaile ibiltari} zaharrak ere
\textsuperscript{(69)}, zeinari atezainak galtzerdiak erosten baitzizkion, atarira
sartu behar izan zuen salmenta amaitzeko.

%%K t14-16-2 p
Erregimentua kuartelera arte lagundu genuen \VUgras{eta}, gure egunarekin oso
pozik, etxera itzuli ginen afaltzera.

%%K t14-tit-8 sub
\VUpk{10.2pt}{40}{Lanbide eta ogibide desberdinak.}
\VUsaut{3.0mm}

%%K t14-17-1 p
17. --- Hurrengo egunean aitak proposatu zigun tokiko egunkariaren inprimategira
eramatea. Pozarren onartu genuen.

%%K t14-17-2 p
\VUgras{Inprimatzailea} \VUgras{liburu-saltzailea} ere bada
\textsuperscript{(70)}, \VUgras{argitaratzailea}, \VUgras{paper-saltzailea} eta
\VUgras{koadernatzailea}. Lehen solairuan jarri ditu egunkariaren
\VUgras{erredakzioa} \textsuperscript{(71)}, eta baita \VUgras{grabatze-tailerra}
ere, non \VUgras{litografoek} \textsuperscript{(72)} eta \VUgras{kobre-grabatzaileek}
\textsuperscript{(73)} lan egiten baitute, eta, azkenik, konposaketa-gela, non
\VUgras{konposatzaileak} \textsuperscript{(74)} eserita baitaude. \VUgras{Makina-aretoa}
\textsuperscript{(75)} bera, \VUgras{inprimatzeko makinak} eta makina desberdinak
beheko solairuan daude.

%%K t14-tit-9 sub
\VUpk{10.2pt}{40}{Joanesek galdera pila bat egiten du.}
\VUsaut{3.0mm}

%%K t14-18-1 p
18. --- Inprimategitik ateratzean, Joanes harridura handiz gelditu zen
\VUgras{mertzeria-dendaren} \textsuperscript{(77)} parean zutoin batean zegoen
beirazko kutxatxo baten aurrean, eta zertarako zen galdetu zuen. Aitak azaldu zion
\VUgras{sute-abisagailua} \textsuperscript{(76)} zela. Eta egia esan, hirian dena
zen mirari bat nire adiskidearentzat: bai \VUgras{edateko iturritxo}
soila \textsuperscript{(78)}, bai \VUgras{banaketa-triziklo} arina
\textsuperscript{(80)}, non \VUgras{mandatari} bat \textsuperscript{(79)} baitzihoan
libreaz jantzita, bai \VUgras{posta-furgoia} \textsuperscript{(81)}.

%%K t14-19-1 p
19. --- Aitak minutu batez utzi gintuen \VUgras{zerga-biltzailearekin}
\textsuperscript{(83)} hitz egiteko --- zeina \VUgras{abokatuarekin}
\textsuperscript{(82)} eta \VUgras{notarioarekin} \textsuperscript{(84)} solasean
gelditu baitzen ---, eta guk kaleko mugimenduari begira olgatu ginen. Jende mota
guztia igarotzen zen gure aurretik: \VUgras{gozotegiko mutila}
\textsuperscript{(85)}, saskian \VUgras{pastel} bikain bat zeramana,
\VUgras{trapu-biltzailearen} \textsuperscript{(86)} atzetik zihoan;
\VUgras{farolzaina} \textsuperscript{(87)} \VUgras{gas-langilearekin}
\textsuperscript{(88)} solasean ari zen; \VUgras{moja} batek
\textsuperscript{(89)}, umezurtz txiki bat eskutik zeramanak, bere komentura
itzultzen zen.

%%K t14-tit-10 sub
\VUpk{10.2pt}{40}{Itzulerako bidean.}
\VUsaut{3.0mm}

%%K t14-20-1 p
20. --- Aita gugana itzuli zen unean bertan, Joanes barrez lehertu zen bi
\VUgras{tximini-garbitzaileren} \textsuperscript{(91)} aurpegi zikinak eta jantzi
bitxia ikusita, kedarrez beltz.

%%K t14-21-1 p
21. --- Amari landare bat erosi nahi nion \VUgras{lore-saltzailearenean}
\textsuperscript{(92)}, baina aitak ohartarazi zuen astun izango zela eramateko eta
hobe zela \VUgras{laranjak} hartzea. \VUgras{Kale-saltzailearengana}
\textsuperscript{(94)} hurbildu ginen, eta \VUgras{gobernantak}
\textsuperscript{(93)}, bere haurrarentzat aukeratzen ari zenak, erosketa amaitu
zuenean, gu zerbitzatu gintuzten.

%%K t14-22-1 p
22. --- Itzulerako bidean ezagun batzuk topatu genituen: gure familiaren
adiskide zahar bat, bere \VUgras{lagunaren} \textsuperscript{(95)} besotik
zihoana, bide eta zubien \VUgras{ingeniari} bat \textsuperscript{(96)} eta nire
lehengusinetako bat bere \VUgras{inudearekin} \textsuperscript{(98)}.

%%K t14-22-2 p
Gero nire amaren \VUgras{apainzalearen} \textsuperscript{(97)} ondotik igaro ginen
eta hirian bisitan zeuden bi \VUgras{fraideren} \textsuperscript{(99)} ondotik,
zeinak aitari geltokirako bidea galdetzera hurbildu baitziren.
\VUsaut{6.0mm}
\VUfilet{22mm}
